\fancyhead[L]{Einleitung} % 1080 Wörter (reste 7920 mots)


\section{Einleitung}


In einer im \cite*{OECD2025GovernmentAtAGlance} erschienen offiziellen Studie kündigte die \textit{Organisation for Economic Co-operation and Development} ein geringes politisches Vertrauen der Bevölkerungen der OECD-Länder an: 
    30\% der Befragten dachten, dass derer Regierung nicht in der Lage sei, dem Einfluss der Unternehmen zu widerstehen, während nur 38\% der Befragten der Meinung waren, dass ihr Parlament die Regierung zur Rechenschaft ziehen könnte.
Laut dieser Studie ist die Wahrnehmung fehlender politische \textit{agency} ein wesentlicher Prädiktor für geringes politisches Vertrauen.
Dieser Gefühl, im eigenen politischen System kein Mitspracherecht zu haben, teilten 30\% der Befragten. 





=> Intérêt du sujet (Actu / Théorique) [150]


-> bp de porto alegre



=> Historique du sujet (Périodisation + 3 auteurs) [230]

    -> 



=> Termes du sujet (Définition, Typologie / Synonymes ou contraire, ce que ce n’Est pas) [200]

    -> différences démocratie délibérative / budget participatif



Dementsprechend sind Bürgerhaushalte sowohl gesellschaftlich als auch wissenschaftlich aktuelle und relevante Forschungsgegenstände.
Im Rahmen dieser Arbeit wird die politische Äußerung der Bürger auf lokaler Ebene untersucht, indem die Themen analysiert werden, die in den Beiträgen zu den Bürgerhaushalten in Stuttgart und Toulouse behandelt werden.
Durch den Vergleich umgesetzter und nicht umgesetzter Vorschläge fügt sich diese Arbeit in eine Reihe wissenschaftlicher Beiträge ein, die die Grenzen des Bürgerhaushalts als Machthebel für die Bevölkerung untersuchen, unter anderem im Rahmen der Untersuchung sogenannter Cherry-Picking-Strategien der institutionellen und politischen Eliten.
Letztendlich ist die vorliegende Arbeit Teil einer umfassenden Forschungsbranche zur deutsch-französischen vergleichenden Politikwissenschaft, da der Fokus dieser Untersuchung auf den beiden Großstädten Toulouse und Stuttgart liegt. \\

Ziel dieser Arbeit ist, die folgende Fragestellung zu beantworten: 
    Inwiefern ähneln oder unterscheiden sich die Themenverteilungen der Bürgerhaushalte in Toulouse und Stuttgart in den Jahren 2021 bis 2023 voneinander?
Zu dieser Forschungsfrage werden die folgenden Unterfragen betrachtet: 
    Gibt es Unterschiede und/oder Ähnlichkeiten bei den Themen, die in den Beiträgen der Bürger zu den Bürgerhaushalten erwähnt werden, und wenn ja, welche?
    Weisen die meist erwähnten Themen in den beiden Städten eine größere Ähnlichkeit auf?
    Sind die Themen der dokumentierten Zusage- und Absagebegründungen in Toulouse und Stuttgart ähnlich ?\\

Im Laufe der vorliegenden Arbeit werden den folgenden Hypothese überprüft:

\begin{itemize}
    \item \textbf{H1:} 
    \item  
\end{itemize}



=> interet de la problematique [250]





=> Methode de ce travail [120]





- Données BP Toulouse et Stuttgart entre 2021 et 2023
- NLP / STM